\begin{proposition}[6.8.2]
\label{IV.2.s6.8.2-fr}
Sean $X$, $Y$ dos preesquemas localmente noetherianos, $f:X\to Y$ un
morfismo y $x$ un punto de $X$. Designemos por $\boldsymbol{M}(f,x)$ una de las propiedades
\textup{(i)} a \textup{(vii)} de la definición \textup{(6.8.1)}, o la propiedad de que $f$ sea suave en el
punto $x$. Sean $Y'$ un preesquema localmente noetheriano, $g:Y'\to Y$ un
morfismo, $X'=X\times_Y Y'$, $f'=f_{(Y')}:X'\to Y'$. Se supone que $f$
o $g$ es localmente de tipo finito. Entonces, para todo $x'\in X'$ situado
sobre $x$, la propiedad $\boldsymbol{M}(f,x)$ implica $\boldsymbol{M}(f',x')$.
\end{proposition}

\begin{proof}
Pongamos $y=f(x)$, $y'=f'(x')$; por transitividad de las fibras (\textbf I, 3.6.4),
se tiene
$f'^{-1}(y')=f^{-1}(y)\otimes_{\mathbf{k}(y)}\boldsymbol{k}(y')$; como o bien $f^{-1}(y)$ es localmente de tipo finito sobre $\boldsymbol{k}(y)$, o bien
$\boldsymbol{k}(y')$ es de tipo finito sobre $\boldsymbol{k}(y)$ (\textbf I, 6.4.11), resulta de (6.7.8) que
las propiedades $\boldsymbol{P}(f^{-1}(y),x)$ y $\boldsymbol{P}(f'^{-1}(y'),x')$ son equivalentes;
además, si $f$ es plano en el punto $x$, $f'$ es plano en el punto $x'$
(2.1.4), lo que demuestra la proposición, teniendo en cuenta (1.4.3, (iii)).
\end{proof}
